\documentclass[11pt,a4paper]{article}
\usepackage[T1]{fontenc}
\usepackage[utf8]{inputenc}
\usepackage{lmodern}
\usepackage[czech]{babel}
\usepackage[margin=1.7cm]{geometry}
\usepackage{amsmath,amssymb,mathtools}
\usepackage{nccmath}
\usepackage{xcolor,cancel,enumitem}
\usepackage{tikz,pgfplots}
\pgfplotsset{compat=1.18}
\usetikzlibrary{arrows.meta,calc,patterns,matrix,shapes.geometric,babel}
\pagestyle{empty}
\setlength{\parindent}{0pt}
\setlength{\parskip}{4pt}
\newcommand{\vysl}[1]{\boxed{\color{red!75!black}#1}}
\newcommand{\pozn}[1]{{\color{red!75!black}#1}}
\newcommand{\dx}{\,\mathrm{d}x}
\newcommand{\dt}{\,\mathrm{d}t}
\newcommand{\subst}[1]{\left|\begin{array}{l}#1\end{array}\right|}
\DeclareMathOperator{\arctg}{arctg}
\DeclareMathOperator{\arccotg}{arccotg}
\DeclareMathOperator{\tg}{tg}
\DeclareMathOperator{\cotg}{cotg}
\usepackage{tcolorbox}
\newtcolorbox{pomocny}{colframe=red!75!black, colback=white, coltitle=white, fonttitle=\bfseries, title={Pomocný výpočet:}, sharp corners}
\newcommand{\dy}{\,\mathrm{d}y}
% znaménkové schéma: \znam{bod/popis/typ,...}{znaménka v intervalech}
% typ: o = vyloučený bod (kroužek), x = obyčejný bod
\newcommand{\znam}[2]{%
\begin{tikzpicture}[baseline=-3pt]
\def\n{0}
\foreach \p/\l/\t [count=\i] in {#1} {\xdef\n{\i}}
\pgfmathsetmacro{\L}{1.6*(\n+1)}
\draw (0,0) -- (\L,0);
\foreach \p/\l/\t [count=\i] in {#1} {
  \pgfmathsetmacro{\xx}{1.6*\i}
  \ifnum\pdfstrcmp{\t}{o}=0 \draw[fill=white] (\xx,0) circle (2.5pt); \draw (\xx-0.09,0) -- (\xx+0.09,0) (\xx,-0.09) -- (\xx,0.09);
  \else \draw (\xx,0.1) -- (\xx,-0.1); \fi
  \node[below=2pt] at (\xx,0) {\small $\l$};
}
\foreach \s [count=\i] in {#2} {
  \pgfmathsetmacro{\xx}{1.6*\i-0.8}
  \node[above] at (\xx,0) {$\s$};
}
\end{tikzpicture}}
\newcommand{\aug}[1]{\left(\begin{array}{rrr|rrr}#1\end{array}\right)}
\newcommand{\krok}[1]{\text{\small\color{orange!80!black}#1}}

\tikzset{
  am/.style={matrix of math nodes, left delimiter=(, right delimiter=),
    inner sep=1.5pt, column sep=0.2cm, row sep=0pt,
    nodes={text width=1.6em, align=right, minimum height=1.45em, inner xsep=0pt, inner ysep=0pt, anchor=center},
    column 1/.style={nodes={text width=0.6em}},
    column 6/.style={nodes={text width=1.1em, align=left, inner xsep=0.03cm}}},
  sip/.style={-{Straight Barb[length=1.6mm]}, semithick}}
% svislá čára rozšířené matice + souřadnice vpravo od matice
\newcommand{\ozn}[1]{\tikz[baseline=(o.base)]\node[draw, ellipse, inner sep=1pt, font=\small] (o) {$#1$};}
\newcommand{\augbar}[1]{%
  \draw ([xshift=-0.1cm]#1-1-6.north west) -- ([xshift=-0.1cm]#1-4-6.south west);
  \coordinate (#1R) at ($(#1.east)+(0.35,0)$);}

\begin{document}
\section*{Soustava, matice -- řešení}

\[
\begin{array}{rcrcrcrcrcr}
3x_1 &-& 2x_2 &+& 6x_3 &+& 2x_4 &-& 4x_5 &=& 5\\
x_1 && &+& 2x_3 &-& x_4 &+& 2x_5 &=& 3\\
x_1 &+& 2x_2 &+& 2x_3 && && &=& 1\\
2x_1 &-& 6x_2 &+& 4x_3 &+& 2x_4 &-& 4x_5 &=& 5
\end{array}
\]
\hrule
\medskip
\begin{center}
\begin{tikzpicture}[baseline]
\matrix (a) [am, anchor=west] at (0,0) {3 & -2 & 6 & 2 & -4 & 5\\ 1 & 0 & 2 & -1 & 2 & 3\\ 1 & 2 & 2 & 0 & 0 & 1\\ 2 & -6 & 4 & 2 & -4 & 5\\};
\augbar{a}
\draw[sip, {Straight Barb[length=1.6mm]}-{Straight Barb[length=1.6mm]}] (aR |- a-1-6) .. controls +(0.7,0) and +(0.7,0) .. (aR |- a-2-6);
\node at ($(aR)+(1.1,0)$) {$\sim$};
\matrix (b) [am, anchor=west] at ($(aR)+(1.6,0)$) {1 & 0 & 2 & -1 & 2 & 3\\ 3 & -2 & 6 & 2 & -4 & 5\\ 1 & 2 & 2 & 0 & 0 & 1\\ 2 & -6 & 4 & 2 & -4 & 5\\};
\augbar{b}
\node[anchor=base west, inner sep=1pt] (b1) at ($(bR |- b-1-6.base)+(0.1,0)$) {$\cdot(-3)$};
\node[anchor=base west, inner sep=1pt] (b2) at ($(b1.base east)+(0.05,0)$) {$\cdot(-1)$};
\node[anchor=base west, inner sep=1pt] (b3) at ($(b2.base east)+(0.05,0)$) {$\cdot(-2)$};
\draw[sip] (b1.south) |- (bR |- b-2-6);
\draw[sip] (b2.south) |- (bR |- b-3-6);
\draw[sip] (b3.south) |- (bR |- b-4-6);
\end{tikzpicture}

\bigskip
\begin{tikzpicture}[baseline]
\node at (-0.55,0) {$\sim$};
\matrix (c) [am, anchor=west] at (0,0) {1 & 0 & 2 & -1 & 2 & 3\\ 0 & -2 & 0 & 5 & -10 & -4\\ 0 & 2 & 0 & 1 & -2 & -2\\ 0 & -6 & 0 & 4 & -8 & -1\\};
\augbar{c}
\draw[sip] (cR |- c-2-6) .. controls +(0.45,0) and +(0.45,0) .. (cR |- c-3-6);
\node at ($(cR |- c-2-6)!0.5!(cR |- c-3-6)+(0.65,0)$) {$+$};
\node[inner sep=1pt, anchor=base west] (c1) at ($(cR |- c-2-6.base)+(1.0,0)$) {$\cdot(-3)$};
\draw[sip] (c1.south) |- (cR |- c-4-6);
\node at ($(cR)+(2.5,0)$) {$\sim$};
\matrix (d) [am, anchor=west] at ($(cR)+(3.0,0)$) {1 & 0 & 2 & -1 & 2 & 3\\ 0 & -2 & 0 & 5 & -10 & -4\\ 0 & 0 & 0 & 6 & -12 & -6\\ 0 & 0 & 0 & -11 & 22 & 11\\};
\augbar{d}
\node[anchor=base west, font=\large] at (dR |- d-3-6.base) {$\cdot\frac16$};
\node[anchor=base west, font=\large] at (dR |- d-4-6.base) {$\cdot\frac1{11}$};
\node at ($(dR)+(1.3,0)$) {$\sim$};
\end{tikzpicture}

\bigskip
\begin{tikzpicture}[baseline]
\node at (-0.55,0) {$\sim$};
\matrix (e) [am, anchor=west] at (0,0) {1 & 0 & 2 & -1 & 2 & 3\\ 0 & -2 & 0 & 5 & -10 & -4\\ 0 & 0 & 0 & 1 & -2 & -1\\ 0 & 0 & 0 & -1 & 2 & 1\\};
\augbar{e}
\draw[sip] (eR |- e-3-6) .. controls +(0.45,0) and +(0.45,0) .. (eR |- e-4-6);
\node at ($(eR |- e-4-6)+(0.65,0)$) {$+$};
\node at ($(eR)+(1.2,0)$) {$\sim$};
\matrix (f) [am, anchor=west] at ($(eR)+(1.8,0)$) {1 & 0 & 2 & -1 & 2 & 3\\ 0 & -2 & 0 & 5 & -10 & -4\\ 0 & 0 & 0 & 1 & -2 & -1\\ 0 & 0 & 0 & 0 & 0 & 0\\};
\augbar{f}
\foreach \j in {1,...,5}
  \node[above=3pt, font=\small] at ([xshift=-0.3em]f-1-\j.north east) {$x_{\j}$};
\draw[rounded corners=4pt]
  ([xshift=-0.7em]f-1-1.north east) -- ([xshift=-0.7em]f-1-1.south east)
  -- ([xshift=-1.4em]f-2-2.north east) -- ([xshift=-1.4em]f-2-2.south east)
  -- ([xshift=-0.85em]f-3-4.north east) -- ([xshift=-0.85em]f-3-4.south east)
  -- ($(fR |- f-3-6.south)+(0.1,0)$);
\node[draw, ellipse, inner sep=1pt, anchor=west, font=\small] at ($(fR |- f-1-6)+(0.15,0)$) {$\ast\ast\ast$};
\node[draw, ellipse, inner sep=1pt, anchor=west, font=\small] at ($(fR |- f-2-6)+(0.15,0)$) {$\ast\ast$};
\node[draw, circle, inner sep=1pt, anchor=west, font=\small] at ($(fR |- f-3-6)+(0.15,0)$) {$\ast$};
\draw[semithick] ($(f.west |- f-4-1.base)+(-0.1,1.2mm)$) -- ($(f.east |- f-4-6.base)+(0.1,1.2mm)$);
\end{tikzpicture}
\end{center}

\begin{center}
\fcolorbox{red!75!black}{white}{\parbox{0.8\textwidth}{
$h(A)=3$, \quad $h(A_r)=3$, \quad $n=5$ \hfill \pozn{\textsc{Frobeniova věta}}\\
$h(A)=h(A_r)$ $\to$ má řešení\\
$h(A)<n$ $\to$ $\infty$ řešení, \quad $n-h(A)=5-3=2$ parametry}}
\end{center}

$\underline{x_5=t}$
\begin{align*}
\ozn{\ast} &\to\ x_4-2x_5=-1 \ \Rightarrow\ \underline{x_4=2x_5-1=2t-1}\\
\ozn{\ast\ast} &\to\ -2x_2+5x_4-10x_5=-4\\
&\phantom{\to\ } -2x_2+5(2t-1)-10t=-4\\
&\phantom{\to\ } -2x_2+10t-5-10t=-4\\
&\phantom{\to\ } -2x_2=-4+5=1 \ \Rightarrow\ \underline{x_2=-\tfrac12}\\
\ozn{\ast\ast\ast} &\to\ x_1+0x_2+2x_3-x_4+2x_5=3\\
&\phantom{\to\ } x_1+2x_3-(2t-1)+2t=3\\
&\phantom{\to\ } x_1+2x_3-2t+1+2t=3\\
&\phantom{\to\ } x_1+2x_3=3-1=2 \ \Rightarrow\ \underline{x_3=u}\\
&\phantom{\to\ } x_1=2-2x_3=\underline{2-2u}
\end{align*}
\[
\vysl{x_1=2-2u,\ \ x_2=-\tfrac12,\ \ x_3=u,\ \ x_4=2t-1,\ \ x_5=t,\quad u,t\in\mathbb{R}}
\]

\newpage
\begin{fleqn}[0pt]
\begin{align*}
A\cdot B &= \begin{pmatrix} 1 & 2 & 4\\ -1 & 0 & 5 \end{pmatrix}\cdot\begin{pmatrix} 1 & 2 & 8\\ 5 & 4 & 6\\ 7 & 9 & 3 \end{pmatrix} =\\
&= \begin{pmatrix} 1\cdot1+2\cdot5+4\cdot7 & 1\cdot2+2\cdot4+4\cdot9 & 1\cdot8+2\cdot6+4\cdot3\\ -1\cdot1+0\cdot5+5\cdot7 & -1\cdot2+0\cdot4+5\cdot9 & -1\cdot8+0\cdot6+5\cdot3 \end{pmatrix}
= \vysl{\begin{pmatrix} 39 & 46 & 32\\ 34 & 43 & 7 \end{pmatrix}}
\end{align*}
\[
\underset{3,3}{B}\cdot\underset{2,3}{A} \ \to\ \text{nelze násobit}
\]
\begin{align*}
A\cdot C &= \begin{pmatrix} 1 & 2 & 4\\ -1 & 0 & 5 \end{pmatrix}\cdot\begin{pmatrix} 5 & 1\\ -1 & 0\\ 2 & 3 \end{pmatrix}
= \begin{pmatrix} 1\cdot5+2\cdot(-1)+4\cdot2 & 1\cdot1+2\cdot0+4\cdot3\\ -1\cdot5+0\cdot(-1)+5\cdot2 & -1\cdot1+0\cdot0+5\cdot3 \end{pmatrix}
= \vysl{\begin{pmatrix} 11 & 13\\ 5 & 14 \end{pmatrix}}
\end{align*}
{\setlength{\arraycolsep}{4.5pt}
\begin{align*}
C\cdot A &= \begin{pmatrix} 5 & 1\\ -1 & 0\\ 2 & 3 \end{pmatrix}\cdot\begin{pmatrix} 1 & 2 & 4\\ -1 & 0 & 5 \end{pmatrix}
= \begin{pmatrix} 5\cdot1+1\cdot(-1) & 5\cdot2+1\cdot0 & 5\cdot4+1\cdot5\\ -1\cdot1+0\cdot(-1) & -1\cdot2+0\cdot0 & -1\cdot4+0\cdot5\\ 2\cdot1+3\cdot(-1) & 2\cdot2+3\cdot0 & 2\cdot4+3\cdot5 \end{pmatrix} = \vysl{\begin{pmatrix} 4 & 10 & 25\\ -1 & -2 & -4\\ -1 & 4 & 23 \end{pmatrix}}
\end{align*}}
\begin{align*}
B\cdot C &= \begin{pmatrix} 1 & 2 & 8\\ 5 & 4 & 6\\ 7 & 9 & 3 \end{pmatrix}\cdot\begin{pmatrix} 5 & 1\\ -1 & 0\\ 2 & 3 \end{pmatrix}
= \begin{pmatrix} 1\cdot5+2\cdot(-1)+8\cdot2 & 1\cdot1+2\cdot0+8\cdot3\\ 5\cdot5+4\cdot(-1)+6\cdot2 & 5\cdot1+4\cdot0+6\cdot3\\ 7\cdot5+9\cdot(-1)+3\cdot2 & 7\cdot1+9\cdot0+3\cdot3 \end{pmatrix} = \vysl{\begin{pmatrix} 19 & 25\\ 33 & 23\\ 32 & 16 \end{pmatrix}}
\end{align*}
\[
\underset{3,2}{C}\cdot\underset{3,3}{B} \ \to\ \text{nelze násobit}
\]
\end{fleqn}

\end{document}
